\section{Effect of jailbreaks on victim model capabilities}\label{sec:capabilities}

\pseudopara{Jailbreak performance on our full dataset.} Here, we evaluate 37 jailbreak methods using our full dataset of 313 forbidden prompts scored by our StrongREJECT rubric-based evaluator on three LLMs of varying capabilities: GPT-4o \cite{GPT4modelcard2024}, GPT-3.5 Turbo \cite{openai_gpt-3.5_turbo}, and Llama-3.1 70B Instruct \cite{Llama31modelcard2024}. All three LLMs are \textit{aligned} in the sense that they were fine-tuned for safety and refuse to respond to our forbidden prompts by default.

The results, shown in \cref{fig:jailbreaks_by_models}, demonstrate that a small number of the jailbreaks we tested are highly effective. The most effective jailbreaks used LLMs to jailbreak LLMs. Prompt Automatic Iterative Refinement (PAIR) instructs an attacker model to iteratively modify a forbidden prompt until it obtains a useful response from the victim model \cite{chao2023jailbreaking}. Persuasive Adversarial Prompts (PAP) instructs an attacker model to persuade a victim model to give it harmful information using techniques like misrepresentation and logical appeals \cite{zeng2024johnny}. Two versions of GCG \cite{zou2023universal, mazeika2024harmbench} are also effective against GPT-3.5 Turbo but not GPT-4o or Llama-3.1 70B, suggesting that newer models may be trained against these adversarial suffixes. Additionally, if an attacker were to iterate through the jailbreaks we tested and select the most effective one for each forbidden prompt (the ``Best" jailbreak in the top row of \cref{fig:jailbreaks_by_models}), they could nearly always achieve a perfect StrongREJECT score, suggesting that frontier models remain highly vulnerable to jailbreaks.

\begin{figure}
    \centering
    \includegraphics[width=0.6\linewidth]{figures/score_by_jailbreak_x_model.pdf}
    \caption{Average StrongREJECT score across various jailbreaks and victim models.}
    \label{fig:jailbreaks_by_models}
\end{figure}

However, most jailbreaks we tested did not result in high-quality responses to forbidden prompts. The jailbreak scores we observe are usually substantially lower than those reported in the papers where these jailbreaks were introduced \cite{wei2023jailbroken}. Combined with our findings in ~\Cref{sec:human_eval}, these results demonstrate that there is a significant discrepancy between StrongREJECT and human judgments on the one hand and most previous automated evaluators on the other. We begin to explain this discrepancy by noting that our StrongREJECT automated evaluator considers both willingness (non-refusal) and capabilities, whereas most previous automated evaluators over-emphasize willingness while ignoring capabilities. This leads us to consider the novel and surprising hypothesis that jailbreaks generally harm model capabilities.

Qualitatively, jailbreaks - especially those that involve obfuscating queries with unusual encodings - often appear to decrease model capabilities. The hypothesis that jailbreaks harm model capabilities also explains the pattern we observe in ~\Cref{tab:condensed_autograder_results}'s first column; most previous automated evaluators are upwards biased because they give high scores to responses where the victim model is \textbf{willing} to respond but \textbf{incapable} of giving a useful answer. Our StrongREJECT evaluator, which considers both willingness and capabilities, gives lower scores.

We want to rigorously test the hypothesis that jailbreaks harm model capabilities. Unfortunately, we cannot test this hypothesis by feeding jailbroken forbidden prompts to aligned victim models, as we did in previous experiments. In this setting, response quality is determined by both willingness and capabilities, whereas we seek experiments to study the effect of jailbreaks on capabilities alone.

To isolate the effect of jailbreaks on capabilities, we must either feed jailbroken forbidden prompts to an \textit{unaligned} model (i.e., a model that willingly responds to forbidden prompts) or feed jailbroken \textit{benign} prompts to an aligned model. We report the results of both of these experiments below.

\pseudopara{Feeding jailbroken forbidden prompts to an \textit{unaligned} model.} We evaluate the same jailbreak methods using our full \ours dataset of 313 forbidden prompts and evaluator on an unaligned model; Dolphin. Because Dolphin is unaligned, jailbreaks will not affect its willingness to respond to forbidden prompts. Therefore, any difference in StrongREJECT scores across jailbreaks must be due to the effect of these jailbreaks on Dolphin’s capabilities.

\cref{fig:willingness_capabilities} (left panel) shows the results of this experiment. The x-axis represents the non-refusal rates (willingness), as measured by string matching, achieved by each of the 37 jailbreaks we tested, averaged across three aligned models (GPT-4o, GPT-3.5 Turbo, and Llama-3.1 70B Instruct) for our full dataset of forbidden prompts. The y-axis represents Dolphin's StrongREJECT scores on the same set of forbidden prompts for each jailbreak. The jailbreaks that most successfully increase aligned models' willingness to respond to forbidden prompts tend to decrease Dolphin's capabilities.

\begin{figure}
    \centering
    \includegraphics[width=\linewidth]{figures/willingness_capabilities.pdf}
    \caption{Jailbreaks that increase an aligned model's willingness to respond to forbidden prompts decrease an unaligned model's ability to respond to forbidden prompts (left) and an aligned model's ability to respond to benign prompts (right). The x-axis represents non-refusal rates averaged across all responses from the aligned models we tested to our full dataset of forbidden prompts. The left-panel y-axis represents StrongREJECT scores for Dolphin's responses to our full dataset of forbidden prompts. The right-panel y-axis represents GPT-4o's MMLU scores. Each point is a jailbreak.}
    \label{fig:willingness_capabilities}
\end{figure}

\pseudopara{Feeding jailbroken \textit{benign} prompts to an aligned model.} The other way to isolate the effect of jailbreaks on model capabilities is to feed jailbroken benign prompts into an aligned model, such as GPT-4o. That is, we start with an underlying benign prompt, modify it according to a jailbreak algorithm, and then feed the resulting ``jailbroken benign prompt” into GPT-4o. For this experiment, we drew benign prompts from the Massive Multitask Language Understanding (MMLU) dataset; a standard test of model capabilities with multiple-choice questions spanning 57 subjects across various disciplines \citep{hendryckstest2021}. Because the underlying prompts are benign, jailbreaks will not affect GPT-4o’s willingness to respond to them. Therefore, any difference in MMLU performance across jailbreaks must be due to the effect of these jailbreaks on GPT-4o’s capabilities.

\cref{fig:willingness_capabilities} (right panel) shows the results of this experiment. The x-axis represents non-refusal rates, as in the left panel. The y-axis represents GPT-4o's zero-shot score on a randomly-selected 171-question subset of MMLU tasks. Similar to the previous experiment, the jailbreaks that most successfully increase aligned models' willingness to respond to forbidden prompts tend to decrease GPT-4o's capabilities as measured by MMLU.

Together, the results of these experiments confirm our hypothesis that jailbreaks generally harm model capabilities, making responses less coherent, less on-topic, less realistic/factual, less detailed, or otherwise lower quality. This finding further underscores the need for researchers to use automated evaluators that, like the StrongREJECT evaluator, consider both willingness and capabilities. In doing so, safety researchers can better direct their efforts towards safeguarding against jailbreaks that, like PAP and PAIR, are genuinely effective.